\chapter{Ataques estructurales: recuperar la clave privada}\label{ch:estructura}

\begin{roadmap}
Los ataques de decodificación del capítulo~\ref{ch:isd} tratan la
clave pública como una matriz aleatoria y van a por el error $e$.
Esta otra familia hace lo contrario: explota que detrás de la matriz
hay \emph{estructura} ---un polinomio de Goppa $g$ y un soporte
ordenado $\alpha_0,\dots,\alpha_{n-1}$--- y ataca directamente la
clave privada. Es la familia que durante casi medio siglo pareció
una promesa rota: soportes que se dividen, distinguidores que
reconocen a Goppa, y aun así un muro de $2^{9116{,}8}$ donde el
mejor ataque genérico cuesta $2^{140{,}8}$. Veremos por qué el diseño
se blindó contra ella ($n<q$), mediremos el abismo con una figura, y
dejaremos el terreno preparado para el terremoto de 2025--2026 del
capítulo~\ref{ch:2026}. Este es el capítulo que responde a la
pregunta «¿y si atacan la clave en vez del mensaje?».
\end{roadmap}

\section{Dos cajas fuertes: la anatomía de la clave privada}

Recordemos qué hay dentro de una clave privada (capítulo~\ref{ch:goppa}):
un polinomio de Goppa
\[
  g(x) \;=\; g_t x^{t} + g_{t-1}x^{t-1} + \cdots + g_1 x + g_0
  \;\in\; \GF{q}[x],
\]
sin raíces en $\GF{q}$, y un soporte ordenado
$\alpha_0,\dots,\alpha_{n-1}$ de elementos distintos de $\GF{q}$
---las raíces del polinomio $\prod_i (x-\alpha_i)$---. Con la clave
privada $(g,\alpha_0,\dots,\alpha_{n-1})$ uno descifra exactamente
igual que el receptor: cualquier atacante que la recupere rompe la
tarea OW-CPA sin tocar el problema SD. La pregunta estructural es,
por tanto: \emph{¿cuánto cuesta enumerar las claves privadas
posibles?}

Dos hechos de conteo elemental, que la guía y la rationale del
equipo subrayan~\cite{cmsecurity,cmrationale}:

\begin{itemize}[leftmargin=1.5em,itemsep=3pt]
  \item Los polinomios $g$ candidatos (mónicos, de grado $t$ sobre
        $\GF{q}$, sin raíces en $\GF{q}$) son del orden de $q^{t}/t$.
        Para el conjunto clásico $m=12$, $t=64$: $q^{t}/t =
        4096^{64}/64 = 2^{768-6} = 2^{762}$. Para los parámetros
        originales de 1978, $2^{494{,}4}$.
  \item Las tuplas ordenadas $(\alpha_0,\dots,\alpha_{n-1})$ de
        elementos distintos de $\GF{q}$ son $q(q-1)\cdots(q-n+1)$:
        para $\texttt{mceliece348864}$, unos $38\,500$ bits
        ($\approx 2^{38\,501}$); para los parámetros originales
        ($n=q$), $1024!\approx 2^{8771}$.
\end{itemize}

A primera vista ambos números son abrumadores, y ahí está el primer
malentendido: \emph{los ataques estructurales no enumeran a lo
bruto}. Dividen el trabajo. La figura~\ref{fig:est-rutas} muestra
las dos rutas conocidas hacia la clave, y la caja negra en la que se
apoyan: el algoritmo de división del soporte de Sendrier
(2000)~\cite{sendrier2000}.

\begin{figure}[htbp]
\centering
\begin{tikzpicture}[
    nd/.style={draw=mcgray,fill=white,rounded corners=3pt,align=flush center,
      font=\scriptsize,inner sep=4pt,minimum height=0.95cm,text width=2.5cm},
    pk/.style={nd,draw=mcblue,fill=mcsky!45,text width=2.2cm,font=\footnotesize},
    ok/.style={nd,draw=mcgreen,fill=mcgreen2!50,text width=2.9cm},
    sup/.style={nd,draw=mcorange,fill=mcorange2!50,text width=3.1cm},
    gen/.style={nd,draw=mcteal,fill=mcteal2!55,text width=2.9cm},
    arr/.style={-{Stealth[length=2.2mm]},thick,mcgray},
    ann/.style={font=\tiny,color=mcgray,align=center}
  ]
  % fila estructural
  \node[pk] (pk1) at (0,2.3) {clave pública\\ ($g$ y $\alpha$'s ocultos)};
  \node[nd,draw=mcorange] (g1) at (3.5,2.3) {adivinar $g$:\\ $2^{762}$ opciones\\ (caso $m=12$, $t=64$)};
  \node[sup] (ss) at (7.3,2.3) {support splitting\\ (caja negra: usa $g$ y el\\ conjunto $\{\alpha_i\}$)};
  \node[ok] (priv) at (11.4,2.3) {$(\alpha_0,\dots,\alpha_{n-1})$ ordenados:\\ descifrar como el receptor};
  \draw[arr] (pk1) -- (g1);
  \draw[arr] (g1) -- (ss);
  \draw[arr] (ss) -- (priv);
  % puerta de abajo
  \node[nd,draw=mcred,fill=mcred2!45] (a1) at (3.5,-1.3)
    {adivinar $(\alpha_0,\dots,\alpha_{n-k})$:\\ $q(q-1)\cdots(q-n+k)$ opciones,\\ $\approx 2^{9116{,}8}$ para \texttt{348864}};
  \draw[arr,mcred] (a1) -- (ss) ;
  \draw[arr,mcred,dashed] (pk1.south) |- (a1.west);
  % fila ISD
  \node[pk] (pk2) at (0,-3.7) {clave pública\\ ($g$ y $\alpha$'s ocultos)};
  \node[gen] (isd) at (3.5,-3.7) {decodificar (ISD,\\ capítulo~\ref{ch:isd}):\\ $2^{140{,}8}$};
  \node[ok] (e2) at (7.3,-3.7) {recuperar $e$:\\ descifrar};
  \draw[arr] (pk2) -- (isd);
  \draw[arr] (isd) -- (e2);
  \draw[arr,dashed,mcgray] (pk2.north) -- (pk1.south)
    node[midway,right,ann,text width=1.6cm]{el atacante elige ruta};
\end{tikzpicture}
\caption{Las rutas conocidas hacia la clave privada, a fecha de 2022.
Arriba, la ruta estructural: adivinar $g$ y dividir el soporte con el
algoritmo de Sendrier~\cite{sendrier2000}; con $n=q$ basta adivinar
$g$, y con $n<q$ la puerta exige además las posiciones
$(\alpha_0,\dots,\alpha_{n-k})$ (ataque de Kirshanova--May, 2022, con
dato parcial~\cite{kirshanova2022}). Abajo, la ruta genérica de
decodificación, que no usa estructura alguna. Los números son los de
las fuentes: $2^{762}$ y $2^{9116{,}8}$ para \texttt{mceliece348864},
$2^{140{,}8}$ de la tabla~\ref{tab:isd-guia}~\cite{cmsecurity}.}
\label{fig:est-rutas}
\end{figure}

\section{La división del soporte}\label{sec:est-support}

El algoritmo de Sendrier merece presentación pausada, porque es la
pieza que hace pensables los ataques estructurales. Su tesis:
\emph{dados $g$ y el desordenado conjunto $\{\alpha_0,\dots,\alpha_{n-1}\}$
---no las posiciones, solo el conjunto---, el polinomio de Goppa
permite reconstruir el orden $\alpha_0,\dots,\alpha_{n-1}$ en tiempo
rápido}. El nombre viene de ahí: el algoritmo «divide» el soporte en
las partes que $g$ distingue, y de esas divisiones deduce qué
columna del código corresponde a qué elemento del
cuerpo~\cite{sendrier2000}. La guía lo enuncia así: recupera
$(\alpha_0,\dots,\alpha_{n-1})$ «a partir de la clave pública, $g$
y $\{\alpha_0,\dots,\alpha_{n-1}\}$, bajo supuestos
menores»~\cite{cmsecurity}. Los detalles finos viven en el artículo
original; para contar el ataque completo basta tratarlo como caja
negra.

La consecuencia para el McEliece original es inmediata. Su diseño
tomaba $n=q$: el soporte es \emph{todo} el cuerpo $\GF{q}$, un
conjunto público, y por tanto el ataque estructural se reduce a
enumerar el polinomio $g$ ---$q^t/t\approx 2^{494{,}4}$ posibilidades
para los parámetros de 1978--- y dejar que la división del soporte
haga el resto. En 2000 esto no rompió a McEliece, pero convirtió lo
que parecía un espacio de claves de $1024!\approx 2^{8771}$ (las
tuplas ordenadas completas) en una enumeración de $q^{t}/t\approx
2^{494{,}4}$ polinomios $g$ ---el orden del soporte lo aporta el
propio algoritmo---, y dejó claro cuál era la pieza a defender.

La defensa, ya lo habrá adivinado el lector, es $n<q$: si el cuerpo
tiene $q$ elementos y solo se usan $n$, el conjunto
$\{\alpha_0,\dots,\alpha_{n-1}\}$ \emph{deja de ser público}.
Aparecen $\binom{q}{n}$ conjuntos posibles, y la caja negra ya no
tiene una de sus entradas. La rationale del equipo lo dice
explícitamente: mantener $n$ por debajo de $q$ es «una defensa extra
contra la división del soporte»; para \texttt{mceliece6960119} hay
$\binom{8192}{6960}\approx 2^{4997}$ conjuntos posibles ---«aunque
podría haber maneras de fusionar el trabajo entre
posibilidades»~\cite{cmrationale}. (Merece la pena notar el doble
filo: contra los ataques de decodificación, $n<q$ no ayuda ni
perjudica; el capítulo~\ref{ch:isd} cuantificó ese efecto en menos
de un bit. La asimetría es real: el mismo truco que es neutro contra
ISD es el que sostiene el muro contra lo estructural.)

En 2022, Kirshanova y May añadieron una pieza más al
rompecabezas~\cite{kirshanova2022}: dado un \emph{dato parcial} ---las
posiciones $(\alpha_0,\dots,\alpha_{n-k})$ de la clave ordenada---,
se recupera el resto de la clave privada con un algoritmo rápido.
Combinado con una búsqueda a lo bruto sobre ese dato parcial, esto
da la cota superior estándar para la recuperación de clave: hay
$q(q-1)\cdots(q-n+k)$ tuplas posibles, «aproximadamente
$2^{9116{,}797}$ posibilidades para
$(q,n,k,t)=(4096,3488,2720,64)$», los parámetros de
\texttt{mceliece348864}~\cite{cmsecurity}. La tabla~\ref{tab:est-espacios}
resume todos los espacios de búsqueda en juego.

\begin{table}[htbp]
\centering
\footnotesize
\setlength{\tabcolsep}{5pt}
\begin{tabular}{llll}
\toprule
ruta & espacio a enumerar & coste (ejemplo) & fuente\\
\midrule
adivinar $g$; soporte conocido ($n=q$) & $q^{t}/t$ &
  $2^{762}$ ($q=4096$, $t=64$) & \cite{cmsecurity}\\
adivinar $g$; con $n<q$ & $\times\,\binom{q}{n}$ &
  $2^{4997}$ ($\binom{8192}{6960}$) & \cite{cmrationale}\\
dato parcial $(\alpha_0,\dots,\alpha_{n-k})$ & $q\cdots(q-n+k)$ &
  $2^{9116{,}8}$ (\texttt{348864}) & \cite{cmsecurity,kirshanova2022}\\
parámetros originales (1978) & $q^{t}/t$ \ y \ $q\cdots(q-n+k)$ &
  $2^{494{,}4}$ \ y \ $2^{4794{,}7}$ & \cite{cmrationale}\\
\midrule
\emph{referencia: decodificación genérica} & --- &
  $2^{140{,}8}$ (\texttt{348864}) & \cite{cmsecurity}\\
\bottomrule
\end{tabular}
\caption{Los espacios de búsqueda de los ataques estructurales, con
los ejemplos que citan las fuentes. La última fila no es un ataque
estructural: es el rival genérico del capítulo~\ref{ch:isd}, incluido
para que la comparación de la figura~\ref{fig:est-abismo} tenga a
mano sus dos extremos.}
\label{tab:est-espacios}
\end{table}

\begin{historia}[Sendrier y la tradición de la equivalencia de códigos]
Los ataques estructurales no nacieron contra McEliece: nacieron de
una pregunta más vieja de la teoría de códigos, la
\emph{equivalencia}: ¿cuándo dos matrices generadoras distintas
describen el mismo código con las posiciones permutadas? Sendrier
llevaba años construyendo herramientas para responderla cuando, en
2000, publicó la división del soporte~\cite{sendrier2000}. Desde
entonces la criptografía basada en códigos la adoptó como parte
estándar del oficio; el capítulo de Overbeck--Sendrier del libro
clásico de criptografía post-cuántica la presenta junto a los ataques
de decodificación como la segunda gran familia~\cite{overbeck2009}.
La lista de trabajos de criptoanálisis que la rationale del equipo
recopila ---treinta publicaciones entre 1981 y 2022, de
Canteaut--Sendrier (1998) a Carrier--Debris-Alazard--Meyer-Hilfiger--Tillich
(2022)--- muestra las dos tradiciones entrelazadas durante cuatro
décadas~\cite{cmrationale,canteautsendrier1998,carrier2022}.
\end{historia}

\begin{observacion}{El doble filo de $n<q$, en una línea}{est-nq}
Para ISD, elegir $n<q$ cuesta a lo sumo fracciones de bit
(observación~\ref{obs:isd-acortamiento}). Para lo estructural,
elegir $n<q$ multiplica el trabajo de la división del soporte por el
número de conjuntos posibles ($\approx 2^{4997}$ para
\texttt{6960119}, según~\cite{cmrationale}). Un mismo parámetro,
dos efectos de signo opuesto: por eso los conjuntos se eligen con el
estimador delante, no a ojo.
\end{observacion}

\section{Distinguidores: saben que es Goppa, ¿y qué?}\label{sec:est-distinguidores}

Hay una familia de resultados estructurales que no recupera la
clave ni decodifica: \emph{distingue}. Un distinguidor es un
algoritmo que, mirando la clave pública, decide si proviene de una
Goppa binaria o de una matriz aleatoria con probabilidad de acierto
apreciable, y con un coste que en los parámetros reales sería
ridículo ($2^{40}$, $2^{60}$\dots). Como en el capítulo~\ref{ch:seguridad}
adelantamos en forma de trampa, un distinguidor viola el supuesto
«pk parece aleatoria», pero no implica por sí solo ni recuperar $e$
ni recuperar la clave. La historia de estos ataques es la historia de
sus \emph{márgenes}: todos funcionan en regímenes de parámetros que
McEliece no usa, y todos mueren mucho antes de llegar a los
conjuntos reales. La tabla~\ref{tab:est-margenes} los recoge.

\begin{table}[htbp]
\centering
\footnotesize
\setlength{\tabcolsep}{6pt}
\begin{tabular}{llll}
\toprule
ataque & criterio & funciona en & McEliece usa\\
\midrule
Faugère et al.\ 2013~\cite{faugere2013} & pocos errores &
  $\le 8$ ($n=1024$);\ $\le 20$ ($n=8192$) & $50$;\ $96$--$128$\\
Carrier et al.\ 2022~\cite{carrier2022} & tasa de código baja &
  $<0{,}3$ & $0{,}7$--$0{,}8$\\
Couvreur et al.\ 2017~\cite{couvreur2017} & grado pequeño &
  $m\le 2$ & $m=12,13$\\
\bottomrule
\end{tabular}
\caption{Los distinguidores conocidos y sus márgenes, según la
rationale del equipo~\cite{cmrationale}. La segunda columna dice qué
región del espacio de parámetros explota cada ataque; la tercera,
hasta dónde llega; la cuarta, dónde viven los conjuntos de Classic
McEliece. El ataque de 2013 es el llamado «distinguidor de alta
tasa» de Faugère--Gauthier-Umaña--Otmani--Perret--Tillich; el de
Couvreur--Otmani--Tillich ataca las variantes «wild».}
\label{tab:est-margenes}
\end{table}

\begin{rigor}[Distinguir no es romper]
La lógica, para que no queden cabos sueltos. (1)~El objetivo
criptográfico del capítulo~\ref{ch:seguridad} es OW-CPA, y el
lema~\ref{lem:seg-ow-sd} lo redujo a SD: recuperar $e$. (2)~Un
distinguidor devuelve un bit ---«¿es Goppa?»--- y nada más; no
produce ni $e$ ni $(g,\alpha_0,\dots)$. (3)~No existe reducción
conocida del primero al segundo: que la clave pública sea
reconocible como Goppa no da, hoy por hoy, ninguna vía directa para
invertir SD más rápido que los ataques genéricos. (4)~El «caso
peor lógico» está identificado desde el propio distinguidor de
2013: que alguien encadene el distinguidor con un ataque de
inversión. Los autores del distinguidor y el equipo de Classic
McEliece han vigilado esa posibilidad durante una década ---y los
márgenes de la tabla~\ref{tab:est-margenes} son precisamente la
distancia que la cadena tendría que cubrir. (5)~Nada de esto es un
teorema de imposibilidad: es un historial. El capítulo~\ref{ch:2026}
contará cómo en 2025 un distinguidor directo ---no una cadena---
cambió el tono de la conversación.
\end{rigor}

\begin{trampa}[«Si su familia cayó, la mía también»]
Cuando una propuesta basada en otra familia de códigos se rompe
---y ha ocurrido muchas veces, como recuerda la
rationale~\cite{cmrationale}---, los autores de los ataques suelen
probarlos contra Goppa binaria. La rationale recoge el resultado
invariable: «estos estudios muestran consistentemente que el sistema
de McEliece está mucho más allá de todos los ataques conocidos»
(cada golpe contra Goppa queda en la situación de la
tabla~\ref{tab:est-margenes}: funcionan en regiones que los
parámetros evitan). La lección no es que Goppa sea invulnerable: es
que cada familia tiene su anatomía, y el criptoanálisis se trasplanta
entre familias con mucha menos facilidad de la que sugieren los
titulares. El capítulo~\ref{ch:2026} mostrará la otra cara del mismo
argumento: cuando el trasplante \emph{sí} prendió.
\end{trampa}

\section{El balance de 2022}\label{sec:est-balance}

Con todo lo anterior, la foto a fecha de la guía de seguridad es
nítida: el mejor ataque de recuperación de clave contra
\texttt{mceliece348864} cuesta $2^{9116{,}8}$ (tabla~\ref{tab:est-espacios})
y el mejor ataque genérico, $2^{140{,}8}$ (tabla~\ref{tab:isd-guia}).
La figura~\ref{fig:est-abismo} pone las tres barras relevantes en una
misma escala logarítmica: decodificar, adivinar $g$, adivinar los
$\alpha$'s.

\begin{figure}[htbp]
\centering
\begin{tikzpicture}[
    bar/.style={draw=mcgray,thick,rounded corners=1.5pt},
    lab/.style={font=\scriptsize,align=left},
    labt/.style={font=\tiny,color=mcgray}
  ]
  \def\sc{0.001383}
  % eje
  \draw[mcgray,thick,-{Stealth[length=2.2mm]}] (0,-0.5) -- (13.2,-0.5);
  \foreach \b in {2000,4000,6000,8000}
    \draw[mcgray] (\sc*\b,-0.62) -- (\sc*\b,-0.38)
      node[below,labt] {$\b$};
  % barras
  \fill[mcsky!70] (0,-0.05) rectangle (\sc*140.8,0.55);
  \draw[bar,mcblue] (0,-0.05) rectangle (\sc*140.8,0.55);
  \node[lab,anchor=west] at (0.30,0.25)
    {decodificar (ISD): $2^{140{,}8}$};
  \fill[mcorange2!80] (0,-1.35) rectangle (\sc*762,-0.75);
  \draw[bar,mcorange] (0,-1.35) rectangle (\sc*762,-0.75);
  \node[lab,anchor=west] at (1.15,-1.05)
    {adivinar $g$ con el soporte conocido: $2^{762}$};
  \fill[mcorange2!80] (0,-2.65) rectangle (\sc*9116.8,-2.05);
  \draw[bar,mcorange] (0,-2.65) rectangle (\sc*9116.8,-2.05);
  \node[lab,anchor=west] at (0.4,-2.35)
    {adivinar $(\alpha_0,\dots,\alpha_{n-k})$: $2^{9116{,}8}$};
  % llave entre barras
  \draw[decorate,decoration={brace,amplitude=3pt},
    mcgray] (\sc*140.8,0.75) -- (\sc*9116.8,0.75);
  \node[lab,text=mcgray] at (6.4,1.05)
    {el abismo: un factor $2^{8976}$};
  % nota
  \node[labt,text width=13cm,align=center] at (6.6,-3.3)
    {cada barra mide $\log_2(\text{coste})$ en la misma escala
     lineal: la primera ocupa aquí $\approx 0{,}2$ cm, y la
     última, casi la figura entera};
\end{tikzpicture}
\caption{El balance de 2022 para \texttt{mceliece348864}, en escala
de exponentes (bits). Decodificar cuesta $2^{140{,}8}$; recuperar la
clave por la vía de $g$ cuesta $2^{762}$ (y esta vía requiere haber
resuelto antes el conjunto, o $n=q$); recuperarla por la vía del
dato parcial cuesta $2^{9116{,}8}$. Datos:
tablas~\ref{tab:est-espacios} y~\ref{tab:isd-guia}.}
\label{fig:est-abismo}
\end{figure}

La lectura del abismo es obligada en este punto del libro. Durante
cuarenta y cinco años, todos los intentos de explotar la estructura
de Goppa se estrellaron contra distancias como esa: la cita de la
rationale es tajante ---«los ataques de recuperación de clave del
estado del arte son enormemente más lentos que la decodificación por
conjuntos de información»~\cite{cmrationale}--- y la conclusión
estratégica del equipo fue doble. Primera: conservar los códigos de
Goppa binarios, «la elección tradicional, conservadora y bien
estudiada»~\cite{cmrationale}, pese a que existen alternativas con
claves más pequeñas. Segunda: documentar cada margen, porque los
márgenes \emph{son} la seguridad en esta familia de ataques: no
existe un teorema que diga «Goppa es rígida»; existe la cuenta de
todos los ataques que se han quedado cortos.

Ese es el muro tal como se veía en 2022. El capítulo~\ref{ch:2026}
cuenta cómo, entre 2025 y 2026, aparecieron las primeras grietas
---primero un distinguidor directo, después ataques de recuperación
de clave--- y cómo reaccionaron el equipo, ISO, BSI y NIST.

\section{Ejercicios}

\begin{enumerate}[leftmargin=1.6em,itemsep=4pt]
  \item \textbf{Aritmética del muro.} Verifique con una calculadora:
        (a)~$4096^{64}/64 = 2^{762}$; (b)~$1024^{50}/50 \approx
        2^{494{,}4}$; (c)~$\log_2\bigl(q(q-1)\cdots(q-769+1)\bigr)
        \approx 9116{,}8$ para $q=4096$ (puede usar la identidad
        $\log_2 P(q,r) = \sum_{j=0}^{r-1}\log_2(q-j)$, una hoja de
        cálculo, o la aproximación de Stirling).
  \item \textbf{Conjuntos y tuplas.} Estime $\log_2\binom{4096}{3488}$
        con la fórmula de Stirling para entropías
        ($\approx 2^{2476}$) y compárela con $2^{9116{,}8}$: ¿por qué
        una va con $n$ y la otra con $n-k$? (pista: la segunda pide
        posiciones \emph{concretas} del soporte ordenado, no solo el
        conjunto).
  \item \textbf{El abismo en potencias.} (a)~Calcule
        $9116{,}8-140{,}8$ y exprese el cociente de costes como
        potencia de $2$. (b)~Compare esa potencia con $2^{256}$
        (el margen de AES-256) y comente en dos líneas por qué el
        titular «recuperar claves de McEliece cuesta $2^{9116}$»
        es, además de correcto en 2022, casi cómico como amenaza
        práctica.
  \item \textbf{Distinguir no es romper.} Use el
        recuadro de rigor: escriba la cadena lógica que convertiría
        un distinguidor en un ataque de inversión de SD, señale el
        eslabón que falta y explique qué papel juegan los márgenes
        de la tabla~\ref{tab:est-margenes}.
  \item \textbf{El doble filo de $n<q$.} Con la
        observación~\ref{obs:isd-acortamiento} (capítulo~\ref{ch:isd})
        y la observación~\ref{obs:est-nq}: explique en un párrafo por
        qué los conjuntos de Classic McEliece pueden tomar $n<q$ sin
        pagar seguridad contra ISD y cobrar, en cambio, ventaja
        contra la división del soporte.
  \item \textbf{Tabla de supervivencia.} Ordene los tres
        distinguidores de la tabla~\ref{tab:est-margenes} por
        «holgura» (distancia relativa entre su límite y los
        parámetros reales) y discuta cuál de los tres le parece el
        candidato más peligroso a escalar. (No hay solución única:
        se trata de razonar sobre los regímenes.)
\end{enumerate}
